% The Critical Line Algorithm as implemented in cvxcla
%
% This paper documents the algorithm implemented in src/cvxcla/cla.py
% (with supporting modules first.py, operators.py and types.py).
%
% Build with:  pdflatex cla.tex   (no bibliography tool required)

\documentclass[11pt]{article}

\usepackage[utf8]{inputenc}
\usepackage[T1]{fontenc}
\usepackage{amsmath, amssymb}
\usepackage{bm}
\usepackage{geometry}
\usepackage{graphicx}
\usepackage{hyperref}
\usepackage{booktabs}
\usepackage{xcolor}
\usepackage[round,authoryear]{natbib}

% The algorithm listing below is typeset with only base LaTeX (a boxed minipage
% plus nested enumerate), so the paper builds without the algorithm/algorithmicx
% packages (algorithm.sty / algpseudocode.sty), which are not present in every
% TeX distribution.

\geometry{margin=1in}
\hypersetup{colorlinks=true, linkcolor=blue!50!black, urlcolor=blue!50!black,
  citecolor=blue!50!black}

% --- notation shortcuts -------------------------------------------------
\newcommand{\R}{\mathbb{R}}
\newcommand{\Sig}{\Sigma}
\newcommand{\F}{\mathcal{F}}        % free set
\newcommand{\B}{\mathcal{B}}        % blocked set
\newcommand{\lb}{\ell}              % lower bound
\newcommand{\ub}{u}                 % upper bound
\newcommand{\trans}{^{\mathsf{T}}}
\newcommand{\inv}{^{-1}}
\DeclareMathOperator*{\argmax}{arg\,max}
\DeclareMathOperator*{\argmin}{arg\,min}

\title{\textbf{The Critical Line Algorithm}\\[0.3em]
\large A parametric active-set method for the entire mean--variance efficient frontier,
as implemented in \texttt{cvxcla}}

\author{cvxcla --- Stanford University Convex Optimization Group}
\date{\today}

\begin{document}
\maketitle

\begin{abstract}
We describe the Critical Line Algorithm (CLA) of Markowitz as implemented in the
\texttt{cvxcla} package. The CLA computes the \emph{entire} mean--variance
efficient frontier of a long-only (or box-constrained) portfolio problem exactly,
rather than sampling it point by point. It exploits the fact that the efficient
frontier is piecewise linear in the weight space: between two consecutive
\emph{turning points} the optimal weights are an affine function $w(\lambda) =
\alpha + \lambda\beta$ of a single scalar parameter $\lambda$ that trades expected
return against variance. The algorithm starts at the maximum-return portfolio
($\lambda=\infty$) and walks the frontier down to the minimum-variance portfolio
($\lambda=0$), changing the status of exactly one asset at each turning point.
Each step amounts to a single solve of a reduced Karush--Kuhn--Tucker (KKT) system,
handled here by block elimination through a small Schur complement. We document the
problem formulation, the optimality conditions, the event logic that detects
turning points, the anti-cycling rule used on degenerate problems, and the
covariance-operator abstraction that lets structured risk models (diagonal plus
low rank) avoid ever materialising a dense $n\times n$ matrix.
\end{abstract}

\tableofcontents

\section{Introduction}

The mean--variance portfolio selection problem of \citet{markowitz1952} asks for the
allocation of capital across $n$ assets that minimises risk for a given level of
expected return. Sweeping the return level produces a one-parameter family of optimal
portfolios whose risk--return image is the \emph{efficient frontier}. The Critical
Line Algorithm (CLA), introduced by \citet{markowitz1956} and developed in
\citet{markowitz1959}, computes this frontier \emph{exactly and in full}: it returns a
finite list of \emph{turning points} between which the frontier is described in closed
form, so no discretisation or repeated re-optimisation is required.

This document describes the implementation in the \texttt{cvxcla} package. The core of
the implementation is the class \texttt{CLA} in \texttt{src/cvxcla/cla.py}; it is
supported by \texttt{first.py} (the starting portfolio), \texttt{operators.py} (the
covariance backend), and \texttt{types.py} (the turning-point and frontier data
structures). The notation below follows the code: $\mu$ is the vector of expected
returns, $\Sig$ the covariance, $A,b$ the linear equality constraints, and $\lb,\ub$ the
box bounds on the weights.

\section{Problem formulation}

\subsection{The parametric quadratic program}

Let $w\in\R^n$ be the vector of portfolio weights, $\mu\in\R^n$ the expected returns,
$\Sig\in\R^{n\times n}$ the (symmetric, positive semidefinite) covariance matrix,
$A\in\R^{m\times n}$ and $b\in\R^m$ the equality constraints, and $\lb,\ub\in\R^n$ the
lower and upper bounds. The implementation traces the frontier through the
\emph{return-parametrised} quadratic program
\begin{equation}
\label{eq:qp}
\begin{aligned}
\min_{w\in\R^n}\quad & \tfrac12\, w\trans \Sig\, w \;-\; \lambda\,\mu\trans w \\
\text{subject to}\quad & A w = b, \\
& \lb \le w \le \ub,
\end{aligned}
\end{equation}
where $\lambda \ge 0$ is a scalar that weights expected return against variance. The
canonical instance has a single equality constraint, the budget constraint
$\mathbf{1}\trans w = 1$ (so $A = \mathbf{1}\trans$ and $b = 1$), but
\eqref{eq:qp} and the implementation allow any $m$ linear equalities.

As $\lambda$ decreases from $+\infty$ to $0$, the optimiser of \eqref{eq:qp} moves from
the maximum expected-return portfolio (subject to the constraints) to the global
minimum-variance portfolio. The locus of these optimisers is exactly the efficient
frontier.

\subsection{Optimality conditions}

At a fixed $\lambda$ the program \eqref{eq:qp} is convex, so the Karush--Kuhn--Tucker
(KKT) conditions are necessary and sufficient. Partition the assets into a
\emph{free} set $\F$ (weights strictly inside their bounds) and a \emph{blocked} set
$\B = \F^{c}$ (weights pinned at a bound, $w_i \in \{\lb_i,\ub_i\}$). For the free
assets the stationarity condition holds with no bound multiplier:
\begin{equation}
\label{eq:kkt}
\Sig_{\F\F}\, w_\F + \Sig_{\F\B}\, w_\B + A_\F\trans \nu \;=\; \lambda\, \mu_\F,
\qquad
A_\F\, w_\F + A_\B\, w_\B \;=\; b,
\end{equation}
where $\nu\in\R^m$ is the multiplier of the equality constraints and the subscripts
select the free/blocked rows and columns. The blocked weights $w_\B$ are known (they
sit on $\lb$ or $\ub$), so \eqref{eq:kkt} is a linear system in the unknowns
$(w_\F,\nu)$.

\subsection{Piecewise-linearity: the affine segment}

The crucial structural fact is that, \emph{while the free/blocked partition is fixed},
the right-hand side of \eqref{eq:kkt} is affine in $\lambda$ and the system matrix is
constant. Hence the solution is affine in $\lambda$:
\begin{equation}
\label{eq:segment}
w(\lambda) \;=\; \alpha \;+\; \lambda\,\beta,
\end{equation}
where $\alpha,\beta\in\R^n$ are constant on the segment (the blocked entries of $\beta$
are zero; the blocked entries of $\alpha$ equal the active bound values). The
implementation calls these vectors \texttt{r\_alpha} and \texttt{r\_beta}. Equation
\eqref{eq:segment} is the \emph{critical line}: a straight line in weight space valid
for an interval of $\lambda$. A \emph{turning point} is a value of $\lambda$ at which
the partition $(\F,\B)$ must change, i.e.\ an endpoint of such an interval.

\section{The first turning point}
\label{sec:first}

The walk starts at $\lambda = +\infty$, where variance is irrelevant and \eqref{eq:qp}
degenerates into ``earn the most expected return that the bounds and budget allow.''
This portfolio is computed in closed form by \texttt{init\_algo}
(\texttt{src/cvxcla/first.py}):

\begin{enumerate}
\item Initialise every weight at its lower bound, $w \gets \lb$.
\item Sort the assets by descending expected return.
\item Walk down that order, raising each asset's weight from $\lb_i$ towards $\ub_i$,
      accumulating capital, until the budget $\sum_i w_i = 1$ is reached (within a
      floating-point tolerance). The last asset touched is only \emph{partially}
      filled --- it absorbs the residual $1 - \sum_i w_i$ --- and is the single
      \emph{free} asset of the first turning point. All others are blocked, either at
      their lower bound (not yet reached) or upper bound (fully filled).
\end{enumerate}

The result is ``the smallest subset of assets with highest return whose upper bounds
sum to at least one,'' which is the maximum-return vertex of the feasible polytope. It
is returned as a \texttt{TurningPoint} with $\lambda = \infty$, a weight vector, and a
boolean \texttt{free} mask. The tolerance $1-10^{-12}$ on the budget test matters: the
increment $1 - \sum_i w_i$ can only reach $1$ up to rounding, and without the slack the
loop would skip past the genuinely interior asset and mark the next (zero-weight) asset
as free.

\section{The turning-point iteration}

From the first turning point the algorithm repeatedly: (i) solves the reduced KKT
system for the current partition to obtain the affine segment \eqref{eq:segment};
(ii) finds the largest $\lambda$ below the current one at which an \emph{event}
occurs; (iii) updates the partition by the single asset responsible for that event;
and (iv) records the new turning point. It stops when $\lambda$ reaches $0$.

\subsection{Solving the reduced KKT system by block elimination}
\label{sec:blockelim}

For the free partition $\F$ the system \eqref{eq:kkt} is the saddle-point (KKT) system
\begin{equation}
\label{eq:saddle}
\begin{bmatrix} \Sig_{\F\F} & A_\F\trans \\[2pt] A_\F & 0 \end{bmatrix}
\begin{bmatrix} w_\F \\ \nu \end{bmatrix}
=
\begin{bmatrix} r_1 \\ r_2 \end{bmatrix}.
\end{equation}
Because the right-hand side splits into a constant part and a $\lambda$-part, the
implementation solves \eqref{eq:saddle} for \emph{two} right-hand sides simultaneously
(the ``$\alpha$ system'' and the ``$\beta$ system''):
\[
r_1^{\alpha} = -\,\Sig_{\F\B}\, w_\B, \quad r_2^{\alpha} = b - A_\B\, w_\B;
\qquad
r_1^{\beta} = \mu_\F, \quad r_2^{\beta} = 0 .
\]
Rather than assembling and factoring the full saddle matrix, the code eliminates the
free block first (\emph{block elimination} / Schur complement). With
\[
Y = \Sig_{\F\F}\inv A_\F\trans, \qquad
z^{\alpha} = \Sig_{\F\F}\inv r_1^{\alpha}, \qquad
z^{\beta} = \Sig_{\F\F}\inv r_1^{\beta},
\]
obtained from a \emph{single} multi-column solve against $\Sig_{\F\F}$, the equality
multipliers come from the $m\times m$ Schur complement
$S = A_\F\,\Sig_{\F\F}\inv A_\F\trans = A_\F Y$:
\begin{equation}
\label{eq:schur}
\nu^{\alpha} = S\inv\!\left(A_\F z^{\alpha} - r_2^{\alpha}\right),
\qquad
\nu^{\beta} = S\inv\!\left(A_\F z^{\beta}\right).
\end{equation}
Back-substitution then gives the free part of the affine segment,
\begin{equation}
\label{eq:freesoln}
\alpha_\F = z^{\alpha} - Y\nu^{\alpha},
\qquad
\beta_\F = z^{\beta} - Y\nu^{\beta},
\end{equation}
while the blocked entries are fixed: $\alpha_\B = w_\B$ and $\beta_\B = 0$. The
dominant cost is the multi-right-hand-side solve against the free block
$\Sig_{\F\F}$; the Schur solve \eqref{eq:schur} is only $m\times m$, and for the
standard budget constraint $m=1$, so it is a scalar division.

\subsection{Reduced gradients for blocked assets}

To decide when a blocked asset should re-enter the free set we need the reduced
gradient of the Lagrangian at the blocked coordinates. The implementation forms two
vectors, again split into a constant and a $\lambda$-linear part:
\begin{align}
\gamma &= \Sig\,\alpha + A\trans \nu^{\alpha}, \label{eq:gamma}\\
\delta &= \Sig\,\beta + A\trans \nu^{\beta} - \mu. \label{eq:delta}
\end{align}
For a blocked asset $i$ the quantity $\gamma_i + \lambda\,\delta_i$ is (proportional to)
its bound multiplier along the segment. While that multiplier keeps the correct sign
the asset is correctly blocked; the $\lambda$ at which it would change sign is exactly
the value at which the asset wants to become free. (The covariance products
$\Sig\alpha$ and $\Sig\beta$ are applied through the covariance operator of
\S\ref{sec:operators}, never as an explicit dense multiply when a structured backend is
in use.)

\subsection{Event detection}
\label{sec:events}

Walking down from the current $\lambda$, the segment \eqref{eq:segment} stays optimal
until one of two things happens. The implementation enumerates four event \emph{types},
two for each direction:

\paragraph{A free asset reaches a bound (free $\to$ blocked).}
A free weight evolves as $w_i(\lambda) = \alpha_i + \lambda\beta_i$. It meets a bound at
\begin{equation}
\label{eq:freeevents}
\lambda_i \;=\;
\begin{cases}
\dfrac{\ub_i - \alpha_i}{\beta_i}, & \beta_i < 0 \quad\text{(heads to the \emph{upper} bound)},\\[1.2em]
\dfrac{\lb_i - \alpha_i}{\beta_i}, & \beta_i > 0 \quad\text{(heads to the \emph{lower} bound)}.
\end{cases}
\end{equation}

\paragraph{A blocked asset's multiplier changes sign (blocked $\to$ free).}
A blocked asset wants to become free at
\begin{equation}
\label{eq:blockedevents}
\lambda_i \;=\; -\,\frac{\gamma_i}{\delta_i},
\end{equation}
admissible only when the asset is at its \emph{upper} bound with $\delta_i < 0$, or at
its \emph{lower} bound with $\delta_i > 0$.

These candidate ratios are assembled into a matrix $L\in\R^{n\times 4}$ (one column per
event type), with entries set to $-\infty$ where the event cannot occur. Two
numerical guards apply:

\begin{itemize}
\item \textbf{Slope floor.} Slopes $\beta_i$ and derivatives $\delta_i$ are tested
against $\sqrt{\varepsilon_{\text{mach}}}$, not the user tolerance. A free weight with a
\emph{tiny but nonzero} slope still crosses a bound given a long enough $\lambda$ range,
so filtering at the coarser tolerance would let weights drift out of bounds; only slopes
at floating-point noise level (below $\sqrt{\varepsilon_{\text{mach}}}$) are discarded,
since the enormous ratios they would produce merely amplify rounding error.

\item \textbf{$\lambda$ window.} The segment is valid only for $\lambda$ at or below the
current value, because the frontier is traced with non-increasing $\lambda$. Candidate
ratios above the current $\lambda$ (by more than the tolerance) are spurious crossings
outside the segment and are set to $-\infty$.
\end{itemize}

The next turning point is at $\lambda^{\star} = \max_{i,t} L_{it}$. If
$\lambda^{\star} < 0$ the frontier is complete.

\subsection{Anti-cycling on degenerate problems}
\label{sec:bland}

On degenerate problems --- tied expected returns, duplicated assets, or several
constraints active at once --- multiple events can occur at the same ratio
$\lambda^{\star}$. A naive ``pick any maximiser'' rule can cycle (the same partitions
recurring without progress). The implementation uses a \textbf{Bland-style rule}: among
all events within the tolerance of $\lambda^{\star}$ it selects the one with the lowest
asset index, and within an asset the lowest event-type column. This makes the choice
deterministic and guarantees termination, resolving degeneracies one asset per
iteration.

The asset and event type selected determine the partition update: event types
``free $\to$ blocked'' clear the asset's \texttt{free} flag, while ``blocked $\to$ free''
set it. Exactly one asset changes status. The new weight vector
$w(\lambda^{\star}) = \alpha + \lambda^{\star}\beta$ is stored as the next turning
point.

\subsection{Termination and the minimum-variance endpoint}

The loop runs while $\lambda > 0$. When no admissible event has a positive ratio the
walk has reached the global minimum-variance portfolio; the algorithm appends a final
turning point at $\lambda = 0$ with weights $\alpha$ (the constant part of the current
segment). A hard iteration cap of $100(n+1)$ guards against non-termination: since each
step fixes the status of at least one asset, a correct trace runs $O(n)$ times, and far
exceeding this signals a failure (e.g.\ cycling), which is raised loudly rather than
allowed to hang. Every stored turning point is validated to satisfy the bounds and the
budget constraint before being accepted.

\subsection{The complete procedure}

Algorithm~1 collects the steps.

\begin{center}
\fbox{\begin{minipage}{0.93\textwidth}
\textbf{Algorithm 1.}~Critical Line Algorithm (as implemented in \texttt{cvxcla}).

\smallskip
\textit{Input:} mean $\mu$, covariance $\Sig$, bounds $\lb,\ub$, equality
constraints $A,b$, tolerance $\tau$.

\smallskip
\begin{enumerate}\setlength{\itemsep}{2pt}
\item Compute $w^{(0)} \gets \texttt{init\_algo}(\mu,\lb,\ub)$, the maximum-return
      vertex (\S\ref{sec:first}); append it with $\lambda \gets \infty$.
\item \textbf{While} $\lambda > 0$:
  \begin{enumerate}\setlength{\itemsep}{2pt}
  \item Read the free/blocked partition $(\F,\B)$ of the last turning point;
        classify $\B$ into \emph{at-upper}/\emph{at-lower} and fix $w_\B$ to the
        active bounds.
  \item Solve the free block $\Sig_{\F\F}$ for $Y, z^{\alpha}, z^{\beta}$
        (multi-RHS, \S\ref{sec:blockelim}).
  \item Form the Schur complement $S = A_\F Y$ and solve \eqref{eq:schur} for
        $\nu^{\alpha},\nu^{\beta}$; back-substitute \eqref{eq:freesoln} to get
        $\alpha,\beta$.
  \item Compute the reduced gradients $\gamma,\delta$ via
        \eqref{eq:gamma}--\eqref{eq:delta}.
  \item Assemble the candidate ratios $L$ from
        \eqref{eq:freeevents}--\eqref{eq:blockedevents}; apply the slope floor and
        the $\lambda$ window. Set $\lambda^{\star} \gets \max L$; \textbf{if}
        $\lambda^{\star} < 0$ \textbf{break}.
  \item Pick the event by the Bland-style rule (\S\ref{sec:bland}) and flip the
        \texttt{free} status of that one asset. Set $\lambda \gets \lambda^{\star}$
        and append the turning point $w \gets \alpha + \lambda^{\star}\beta$.
  \end{enumerate}
\item Append the final turning point $w \gets \alpha$ at $\lambda = 0$ (the
      minimum-variance portfolio) and \textbf{return} the list of turning points.
\end{enumerate}
\end{minipage}}
\end{center}

\section{The covariance operator abstraction}
\label{sec:operators}

The turning-point loop touches the covariance matrix through only three operations:
a full matrix--vector product $\Sig x$ (\texttt{matvec}); a solve against the free block,
$\Sig_{\F\F}\inv R$ for a multi-column right-hand side $R$ (\texttt{solve\_free}); and a
free-to-blocked cross product $\Sig_{\F\B}\,x_\B$ (\texttt{cross}). The package captures
exactly this contract in the \texttt{CovarianceOperator} protocol
(\texttt{src/cvxcla/operators.py}), so the algorithm never assumes an explicit matrix.
Two backends implement it.

\subsection{Dense covariance}

\texttt{DenseCovariance} wraps an explicit symmetric $n\times n$ matrix and implements
the three operations with plain \texttt{numpy} calls: \texttt{matvec} is $\Sig x$,
\texttt{solve\_free} is \texttt{np.linalg.solve} on the principal submatrix
$\Sig_{\F\F}$, and \texttt{cross} is the off-diagonal block product. This reproduces
exactly the behaviour of the classical CLA on a raw matrix and is the reference
backend.

\subsection{Factor (diagonal-plus-low-rank) covariance}

\texttt{FactorCovariance} represents a structured risk model
\begin{equation}
\label{eq:factor}
\Sig \;=\; \operatorname{diag}(d) \;+\; U\,\Delta\,U\trans,
\qquad d\in\R^{n}_{>0},\; U\in\R^{n\times k},\; \Delta\in\R^{k\times k},
\end{equation}
the form taken by factor models and by random-matrix-theory-cleaned covariances
(constant-residual-eigenvalue / eigenvalue clipping). The free-block solve never forms
$\Sig_{\F\F}$; instead it uses the Woodbury identity
\begin{equation}
\label{eq:woodbury}
\Sig_{\F\F}\inv \;=\; D_\F\inv \;-\; D_\F\inv U_\F\, W\inv\, U_\F\trans D_\F\inv,
\qquad
W \;=\; \Delta\inv + U_\F\trans D_\F\inv U_\F \;\in\;\R^{k\times k},
\end{equation}
so a solve costs $O(n_\F k^2 + k^3)$ rather than $O(n_\F^3)$, and no $n\times n$ matrix
is ever allocated --- memory and per-solve cost scale as $O(nk)$ instead of $O(n^2)$.
The cross product \eqref{eq:factor} contributes only through its low-rank part, since
the diagonal has no off-diagonal block:
\(
\Sig_{\F\B}\,x_\B = U_\F\,\Delta\,(U_\B\trans x_\B).
\)
Because the CLA accesses the covariance solely through the protocol, switching from a
dense matrix to a factor model requires no change to the turning-point loop.

\section{The efficient frontier object}

The list of turning points is wrapped in a \texttt{Frontier} object
(\texttt{src/cvxcla/types.py}). Each \texttt{FrontierPoint} carries a weight vector and
exposes its expected return $\mu\trans w$ and variance $w\trans\Sig w$. Between adjacent
turning points the frontier in \emph{weight} space is the straight segment
\eqref{eq:segment}, so the object can \texttt{interpolate} any number of intermediate
portfolios by convex combination. Derived quantities --- volatility, Sharpe ratio --- are
computed pointwise. The maximum-Sharpe portfolio is found by a one-dimensional search
over the convex combination of the two turning points bracketing the discrete Sharpe
maximum, exploiting the same piecewise-linear weight structure.

\section{Complexity and properties}

\begin{itemize}
\item \textbf{Exactness.} The frontier is returned exactly as a finite set of turning
points; between them the closed-form segment \eqref{eq:segment} holds. No
discretisation error is introduced.

\item \textbf{Number of turning points.} Each iteration changes exactly one asset's
status, so the number of turning points is $O(n)$ in practice; the implementation caps
iterations at $100(n+1)$ as a safety net.

\item \textbf{Per-iteration cost.} Dominated by the free-block solve: $O(n_\F^3)$ for
the dense backend, or $O(n_\F k^2 + k^3)$ for the factor backend via Woodbury
\eqref{eq:woodbury}. The Schur complement \eqref{eq:schur} adds only an $m\times m$
solve ($m=1$ for the budget constraint).

\item \textbf{Robustness on degenerate data.} The Bland-style tie-break
(\S\ref{sec:bland}) and the $\sqrt{\varepsilon_{\text{mach}}}$ slope floor
(\S\ref{sec:events}) make the trace deterministic and prevent both cycling and
out-of-bounds drift.
\end{itemize}

\section{Numerical experiment}
\label{sec:experiment}

To illustrate the algorithm at a realistic scale we trace the entire efficient
frontier of a $200$-asset problem and then measure how the runtime grows with the
number of assets. Both experiments are reproducible from fixed random seeds:
\texttt{experiments/frontier\_200x2000.py} and
\texttt{experiments/runtime\_scaling.py}.

\subsection{Tracing a 200-asset frontier}

\paragraph{Setup.}
We simulate $T = 2000$ daily returns for $N = 200$ assets from a $K = 15$ factor
model. As the risk model we use the factor covariance itself,
$\Sig = \operatorname{diag}(d) + U\,\Delta\,U\trans$ (standard practice --- a
structured risk model rather than the noisy sample covariance), and we estimate
the expected returns $\mu \in \R^{200}$ by the sample mean over the $2000$ days.
We trace the long-only, fully-invested frontier, i.e.\ \eqref{eq:qp} with the
budget constraint $\mathbf{1}\trans w = 1$ and box bounds $0 \le w \le 1$. The
same frontier is traced twice with two mathematically identical representations
of this $\Sig$: the dense matrix and the structured \texttt{FactorCovariance}
(Woodbury) backend (its condition number here is $\approx 1.0\times 10^{3}$).

\paragraph{Results.}
The algorithm returns the \emph{entire} frontier as a sequence of
\textbf{207 turning points} (Figure~\ref{fig:frontier}); between consecutive
points the frontier is the exact affine segment \eqref{eq:segment}, so these
$207$ corner portfolios describe the whole efficient set with no discretisation.
Because the dense matrix and the factor operator represent the \emph{same} $\Sig$,
the two backends return the identical frontier ($207$ turning points either way) ---
the Woodbury identity is exact, so it changes only the cost of each solve, never
the result. Tracing the frontier takes a median of $\approx 38$~ms with the dense
backend ($\approx 0.18$~ms per turning point) and $\approx 19$~ms with the factor
backend, a $2.0\times$ speedup at this size; Table~\ref{tab:experiment} collects
the figures. (As an aside, replacing $\Sig$ by a low-rank \emph{approximation} of
the full sample covariance --- its leading eigenpairs plus a diagonal residual ---
yields a near-identical frontier differing by at most one corner, which is why
factor risk models are a practical substitute for the raw sample covariance.)

\begin{table}[h]
\centering
\begin{tabular}{lrr}
\toprule
 & \textbf{Dense backend} & \textbf{Factor backend (Woodbury)} \\
\midrule
Turning points         & $207$             & $207$ \\
Median trace time      & $\approx 38$~ms   & $\approx 19$~ms \\
Time per turning point & $\approx 0.18$~ms & $\approx 0.09$~ms \\
\bottomrule
\end{tabular}
\caption{Tracing the full $200$-asset frontier (median of $5$ repetitions). Both
backends solve the identical covariance and return the same $207$-point frontier;
the factor backend never materialises the $200\times 200$ matrix and is about
twice as fast.}
\label{tab:experiment}
\end{table}

\begin{figure}[h]
\centering
\includegraphics[width=0.72\textwidth]{frontier.pdf}
\caption{The efficient frontier of the $200$-asset, $2000$-day problem in
expected-return / volatility space, computed exactly by the Critical Line
Algorithm. Each marker is one of the $207$ turning points; the endpoints are the
maximum-return portfolio ($\lambda=\infty$) and the global minimum-variance
portfolio ($\lambda=0$).}
\label{fig:frontier}
\end{figure}

\subsection{Runtime versus problem size}
\label{sec:scaling}

We next trace the full frontier for $n \in \{20, 40, 80, 160, 320, 640\}$ assets,
each from a ground-truth $K = 10$ factor model so that the dense and factor
backends again solve the identical $\Sig$ and return the same frontier. For
reference we also time PyPortfolioOpt's \texttt{CLA} \citep{bailey2013}, the
widely-used reference implementation of the critical-line algorithm, on the same
dense problem. The number of turning points tracks $n$ closely
($21, 42, 81, 159, 322, 641$ respectively, and PyPortfolioOpt finds essentially
the same set), confirming the $O(n)$ count of \S\ref{sec:experiment}.

Figure~\ref{fig:scaling} plots the median trace time against $n$ on log--log
axes. A least-squares fit over the largest four sizes gives an empirical exponent
of $\approx 2.2$ for the \texttt{cvxcla} dense backend and $\approx 1.3$ for the
factor backend: the dense per-iteration solve scales like $O(n_\F^3)$ over $O(n)$
iterations, while the Woodbury solve with fixed $K$ stays close to linear in $n$.
The dense and factor backends are comparable for small universes (the Woodbury
correction carries a fixed $K\times K$ overhead), but the factor backend pulls
away as $n$ grows --- from rough parity at $n=40$ to a $\approx 8\times$ speedup
at $n=640$ (Table~\ref{tab:scaling}). PyPortfolioOpt's pure-Python implementation
scales far more steeply ($\approx n^{3.4}$): at $n=640$ it takes about
$253$~seconds, so \texttt{cvxcla}'s dense backend is roughly $330\times$ faster
and the factor backend roughly $2800\times$ faster. Much of that gap is
implementation (a vectorised block-elimination over \texttt{numpy} versus
element-wise Python loops); the factor backend adds the asymptotic Woodbury
advantage on top.

\begin{table}[h]
\centering
\begin{tabular}{rrrrr}
\toprule
$n$ & Turning points & PyPortfolioOpt [ms] & \texttt{cvxcla} dense [ms] & \texttt{cvxcla} factor [ms] \\
\midrule
$20$  & $21$  & $7.4$        & $1.3$   & $1.5$ \\
$40$  & $42$  & $40.3$       & $3.2$   & $3.2$ \\
$80$  & $81$  & $212.1$      & $7.2$   & $6.3$ \\
$160$ & $159$ & $1{,}477$    & $23.5$  & $14.3$ \\
$320$ & $322$ & $15{,}932$   & $113.8$ & $33.4$ \\
$640$ & $641$ & $253{,}307$  & $762.7$ & $91.6$ \\
\bottomrule
\end{tabular}
\caption{Frontier trace time versus number of assets $n$ ($K=10$ factors;
\texttt{cvxcla} times are the median of $3$ repetitions, PyPortfolioOpt a single
run). All implementations return the same $O(n)$-point frontier. Empirically the
times scale like $\approx n^{3.4}$ (PyPortfolioOpt), $\approx n^{2.2}$
(\texttt{cvxcla} dense) and $\approx n^{1.3}$ (\texttt{cvxcla} factor).}
\label{tab:scaling}
\end{table}

\begin{figure}[h]
\centering
\includegraphics[width=0.72\textwidth]{scaling.pdf}
\caption{Frontier trace time as a function of the number of assets $n$ (log--log).
PyPortfolioOpt's reference critical-line implementation \citep{bailey2013} scales
like $\approx n^{3.4}$; \texttt{cvxcla}'s dense backend like $\approx n^{2.2}$;
and its diagonal-plus-low-rank factor backend, solving the identical covariance
through the Woodbury identity, like $\approx n^{1.3}$. At $n = 640$ the factor
backend is roughly $2800\times$ faster than PyPortfolioOpt and $8\times$ faster
than the dense backend.}
\label{fig:scaling}
\end{figure}

\paragraph{Remark on conditioning.}
These traces run cleanly because the expected returns are well separated and the
covariance is well conditioned. On strongly degenerate inputs --- (near-)tied
expected returns, duplicated assets, or an ill-conditioned free block --- the
event logic can encounter many coincident ratios; the Bland-style rule of
\S\ref{sec:bland} resolves these deterministically, but very tie-heavy problems
remain numerically delicate.

\section{Conclusion}

The Critical Line Algorithm computes the complete mean--variance efficient frontier by
recognising its piecewise-linear structure in weight space and walking from the
maximum-return portfolio to the minimum-variance portfolio, flipping one asset between
the free and blocked sets at each turning point. The \texttt{cvxcla} implementation
realises each step as a single block-eliminated KKT solve, hardens the event logic
against degeneracy and floating-point noise, and abstracts the covariance behind a
small operator protocol so that structured risk models obtain the same exact frontier
at a fraction of the dense cost.

\begin{thebibliography}{9}

\bibitem[Markowitz(1952)]{markowitz1952}
H.~M. Markowitz.
\newblock Portfolio selection.
\newblock \emph{The Journal of Finance}, 7(1):77--91, 1952.

\bibitem[Markowitz(1956)]{markowitz1956}
H.~M. Markowitz.
\newblock The optimization of a quadratic function subject to linear constraints.
\newblock \emph{Naval Research Logistics Quarterly}, 3(1--2):111--133, 1956.

\bibitem[Markowitz(1959)]{markowitz1959}
H.~M. Markowitz.
\newblock \emph{Portfolio Selection: Efficient Diversification of Investments}.
\newblock John Wiley \& Sons, New York, 1959.

\bibitem[Niedermayer and Niedermayer(2010)]{niedermayer2010}
A.~Niedermayer and D.~Niedermayer.
\newblock Applying Markowitz's critical line algorithm.
\newblock In \emph{Handbook of Portfolio Construction}, pages 383--400. Springer, 2010.

\bibitem[Bailey and L\'opez de Prado(2013)]{bailey2013}
D.~H. Bailey and M.~L\'opez de Prado.
\newblock An open-source implementation of the critical-line algorithm for portfolio
optimization.
\newblock \emph{Algorithms}, 6(1):169--196, 2013.

\end{thebibliography}

\end{document}
